\section{Experiments and Results}%
\label{sec:experiments}

% Plot style definition
\pgfplotsset{
    performance/.style={
        width=\linewidth,
        height=0.6\linewidth,
        grid=major,
        grid style={line width=0.2pt, draw=gray!15},
        axis lines=left,
        xlabel=Agent Iteration,
        ylabel=Utility,
        xlabel near ticks,
        ylabel near ticks,
        legend pos=south east,
        legend cell align=left,
        legend style={
            draw=none,
            fill=white,
            fill opacity=0.8,
            text opacity=1,
            font=\small
        },
        try min ticks=5,
        cycle list={
            {blue!40, opacity=0.4},
            {red!80!black, thick},
            {blue!20},
        },
        tick style={color=gray!70},
        major grid style={line width=0.2pt,draw=gray!15},
        minor grid style={line width=0.1pt,draw=gray!8},
        every axis label/.append style={font=\small},
        tick label style={font=\small}
    }
}

% Performance plot
\begin{figure}[t]
    \centering
    \begin{tikzpicture}
    \begin{axis}[performance]
        % Load data
        \pgfplotstableread{data/main_plot_data.dat}\perfdata

        % Raw scores
        \addplot[c3!90] table[x=iteration,y=target_score] {\perfdata};
        \addlegendentry{Mean Utility}

        % Cumulative maximum
        \addplot table[x=iteration,y=cummax_score] {\perfdata};
        \addlegendentry{Best Utility So Far}

        % Confidence intervals
        \addplot[name path=upper,draw=none] table[x=iteration,y=ci_upper] {\perfdata};
        \addplot[name path=lower,draw=none] table[x=iteration,y=ci_lower] {\perfdata};
        \addplot[c3!30] fill between[of=upper and lower];
        \addlegendentry{95\% Confidence Interval}

        \node[anchor=south, fill=white, fill opacity=0.9, text opacity=1,
              rounded corners=2pt, inner sep=2pt, font=\footnotesize, align=center]
              at (axis cs:1.1,0.63) {`Smart Edit'\\Tool};
        \draw[->,>=stealth,thick,opacity=0.6] (axis cs:1.0,0.63) -- (axis cs:1,0.59);

        \node[anchor=center, fill=white, fill opacity=0.9, text opacity=1,
              rounded corners=2pt, inner sep=2pt, font=\footnotesize, align=center]
              at (axis cs:5.0,0.55) {Code Context\\Summarization};
        \draw[->,>=stealth,thick,opacity=0.6] (axis cs:5.0,0.562) -- (axis cs:5,0.59);

        \node[anchor=center, fill=white, fill opacity=0.9, text opacity=1,
              rounded corners=2pt, inner sep=2pt, font=\footnotesize, align=center]
              at (axis cs:7,0.655) {File Edit\\Verification};
        \draw[->,>=stealth,thick,opacity=0.6] (axis cs:7,0.645) -- (axis cs:7,0.605);

        \node[anchor=center, fill=white, fill opacity=0.9, text opacity=1,
              rounded corners=2pt, inner sep=2pt, font=\footnotesize, align=center]
              at (axis cs:9,0.685) {AST Symbol\\Locator};
        \draw[->,>=stealth,thick,opacity=0.6] (axis cs:9,0.675) -- (axis cs:9,0.65);

        \node[anchor=center, fill=white, fill opacity=0.9, text opacity=1,
              rounded corners=2pt, inner sep=2pt, font=\footnotesize]
              at (axis cs:13,0.7) {Hybrid Symbol Locator};
        \draw[->,>=stealth,thick,opacity=0.6] (axis cs:13,0.693) -- (axis cs:13.9,0.675);
    \end{axis}
    \end{tikzpicture}
    \caption{Performance across iterations. Key improvements are annotated with
      their corresponding tool or agent modifications.}
    \label{fig:performance}
\end{figure}

We first show a run that includes relatively long-horizon tasks in the benchmark
set. We include \acr{SWE} Bench Verified
\citep{chowdhuryIntroducingSWEbenchVerified2024} which tests problem
decomposition, code navigation as well as fast and efficient file editing from
which we select a fixed random subset of 50 questions. We also include 50 random
questions from LiveCodeBench~\citep{jain2024livecodebench}, the questions from
which resemble competitive programming and often include more theoretical
reasoning\footnote{Both SWEBench-Verified and LiveCodeBench are MIT licensed}.

We also include two example `synthetic' benchmark tasks, defined and curated
entirely within the agent's codebase. The first of these is a file editing
benchmark, constructed by cloning the repositories used in \acr{SWE} bench,
finding `interesting' (i.e.~not trivial, nor complete overwrites) file edits in
consecutive commits, mounting the repository checked out at the first commit in
the working directory, and prompting the agent to edit the file such that it
matches the target content of the second commit. The agent is scored on the
closeness of the final file content to the target content, with time and dollar
cost and token consumption also being accounted for in the utility function.
The second is a codebase navigation problem, which is also constructed by
cloning some Python repositories, from which we identify `interesting' symbols,
and at least one reference to this symbol in the project. The benchmark task is
for the agent to locate the definition of this symbol, and return the location
in the format \texttt{path/to/file.py:line\_num:column\_num} from the example
reference.

We run the agent system using Sonnet 3.5 (v2)
\citep{anthropicIntroducingClaude35} in most of the agents in the agent system,
with the exception of a `reasoning agent' which uses o3 mini
\citep{openaiOpenAIO3mini2025}. While the base agent system can run with much
smaller and less expensive models, and many of our benchmark tasks can be solved
by this class of models, the
meta-improvement task is a complicated long-horizon task which requires powerful
models. The API cost of the 15 iteration run was approximately \$7,000.

The results are shown in Figure~\ref{fig:performance}, with more detailed
iteration-level peformance numbers in Table~\ref{tab:performance}.

\begin{table}[t]
\caption{Performance per iteration: the benchmark accuracy numbers correspond to
a synthetic file editing benchmark and symbol location (codebase navigation)
benchmark, SWE-Bench Verified, and LiveCodeBench, respectively. We also show the
metrics averaged across all problem instances across all benchmarks to compare
the relative `per problem' costs across agent iterations.}
\label{tab:performance}
\begin{center}

% Load data
\pgfplotstableread[col sep=space]{data/main_table.dat}\loadeddata

\pgfplotstabletypeset[
    columns={iteration,file,symbol,swe,lcb,cost,time,tokens,cached},
    every head row/.style={
        before row={
            \toprule
            & \multicolumn{4}{c}{\bf Benchmark Accuracy} &
            \multicolumn{4}{c}{\textbf{Average Metrics} (per benchmark problem)} \\
            \cmidrule(lr){2-5} \cmidrule(lr){6-9}
        },
        after row=\midrule
    },
    every last row/.style={after row=\bottomrule},
    columns/iteration/.style={
        column type=r,
        string type,
        column name={\bf Iter}
    },
    columns/file/.style={
        column type=r,
        fixed zerofill,
        precision=2,
        column name={\bf File}
    },
    columns/symbol/.style={
        column type=r,
        fixed zerofill,
        precision=2,
        column name={\bf Sym}
    },
    columns/swe/.style={
        column type=r,
        fixed zerofill,
        precision=2,
        column name={\bf SWE-Bv}
    },
    columns/lcb/.style={
        column type=r,
        fixed zerofill,
        precision=2,
        column name={\bf LCB}
    },
    columns/cost/.style={
        column type=r,
        fixed zerofill,
        precision=2,
        column name={\bf Cost (\$)}
    },
    columns/time/.style={
        column type=r,
        fixed zerofill,
        precision=1,
        column name={\bf Time (s)}
    },
    columns/tokens/.style={
        column type=r,
        string type,
        column name={\bf Tokens (M)}
    },
    columns/cached/.style={
        column type=l,
        string type,
        column name={\bf \% Cached}
    }
]\loadeddata

\end{center}
\end{table}

The numbers show a promising progression. For the file editing benchmark, the
agent is initially only equipped with the ability to overwrite files which is
slow and token inefficient, and it quickly implemented self-termed `smart'
editing tools based on diffs and ranges. For the symbol navigation task, the
agent implemented a simple \acr{AST} based symbol locator at iteration 9
(perhaps exploiting the fact that our benchmarks were predominantly Python
based), which yielded good improvements also reflected in other
tasks.\footnote{Note that we attribute the low score in the symbol locator
  benchmark to poor data quality---some of the target symbols point to
  uninstalled external libraries for instance.} The \acr{SWE} Bench Verified
subset saw an appreciable increase in performance throughout the run, accuring
agent framework improvements. The LiveCodeBench scores saw a subtle
improvement, although not as pronounced as some of the other benchmarks.

\subsection{Performance on Reasoning Tasks}%
\label{subsec:reasoning_tasks}

We also evaluate the effectiveness of the self-referential agent system in
improving task performance in more reasoning-heavy domains. We ran another
experiment with two question answering tasks in the benchmark set: \acr{AIME}
2024 and \acr{GPQA} Diamond \cite{reinGPQAGraduateLevelGoogleProof2023}, the
results of which are shown in Figure~\ref{fig:aime_gpqa}.

\pgfplotsset{
    aime/.style={
        width=\linewidth,
        height=0.4\linewidth,
        grid=major,
        grid style={line width=0.2pt, draw=gray!15},
        axis lines=left,
        title=\acr{AIME} and \acr{GPQA} Show Little Improvement,
        xlabel=Agent Iteration,
        ylabel=Mean Accuracy,
        xlabel near ticks,
        ylabel near ticks,
        legend pos=north east,
        legend cell align=left,
        legend style={
            draw=none,
            fill=white,
            fill opacity=0.8,
            text opacity=1,
            font=\small
        },
        try min ticks=5,
        cycle list={
            {blue!40, opacity=0.4},
            {red!80!black, thick},
            {blue!20},
        },
        tick style={color=gray!70},
        major grid style={line width=0.2pt,draw=gray!25},
        minor grid style={line width=0.1pt,draw=gray!18},
        every axis label/.append style={font=\small},
        tick label style={font=\small}
    }
}

\begin{figure}[htbp]
    \centering
    \begin{tikzpicture}
    \begin{axis}[aime]
        % Load data
        \pgfplotstableread{data/aime_gpqa_plot.dat}\perfdata

        % Raw scores
        \addplot[c3!90] table[x=iteration,y=target_score] {\perfdata};
        \addlegendentry{Mean Accuracy}

        % Cumulative maximum
        \addplot table[x=iteration,y=cummax_score] {\perfdata};
        \addlegendentry{Best Mean Accuracy So Far}

        % Confidence intervals
        \addplot[name path=upper,draw=none] table[x=iteration,y=ci_upper] {\perfdata};
        \addplot[name path=lower,draw=none] table[x=iteration,y=ci_lower] {\perfdata};
        \addplot[c3!30] fill between[of=upper and lower];
        \addlegendentry{95\% Confidence Interval}

        % Annotations for key iterations (adjust coordinates based on your data)
        \node[anchor=south, fill=white, fill opacity=0.9, text opacity=1,
              rounded corners=2pt, inner sep=2pt, font=\footnotesize, align=center]
              at (axis cs:1,0.79) {Independent math\\verifier};
        \draw[->,>=stealth,thick,opacity=0.6] (axis cs:1,0.785) -- (axis cs:1,0.76);

        \node[anchor=north, fill=white, fill opacity=0.9, text opacity=1,
              rounded corners=2pt, inner sep=2pt, font=\footnotesize, align=center]
              at (axis cs:2,0.71) {Sympy symbolic\\calculator};
        \draw[->,>=stealth,thick,opacity=0.6] (axis cs:2,0.715) -- (axis cs:2,0.76);

        \node[anchor=north, fill=white, fill opacity=0.9, text opacity=1,
              rounded corners=2pt, inner sep=2pt, font=\footnotesize, align=center]
              at (axis cs:4,0.66) {Math cross\\validator};
        \draw[->,>=stealth,thick,opacity=0.6] (axis cs:4,0.66) -- (axis cs:4,0.685);

        \node[anchor=south, fill=white, fill opacity=0.9, text opacity=1,
              rounded corners=2pt, inner sep=2pt, font=\footnotesize, align=center]
              at (axis cs:4.8,0.79) {Geometry\\specialist};
        \draw[->,>=stealth,thick,opacity=0.6] (axis cs:5,0.785) -- (axis cs:5,0.77);

        \node[anchor=north, fill=white, fill opacity=0.9, text opacity=1,
              rounded corners=2pt, inner sep=2pt, font=\footnotesize, align=center]
              at (axis cs:6,0.66) {Systematic mathematical\\reasoning};
        \draw[->,>=stealth,thick,opacity=0.6] (axis cs:6,0.66) -- (axis cs:6,0.69);
    \end{axis}
    \end{tikzpicture}
    \caption{Agent Framework Saturation: the benefits the agent system was able
      to find when the models alone (e.g. o3-mini-high) already perform well was
      marginal.}
    \label{fig:aime_gpqa}
\end{figure}

The results here show less improvement, and highlight an important interplay
between the base models and the scaffolding system. The base agent system at
iteration 0 used in Figure~\ref{fig:aime_gpqa} used Sonnet 3.5, with a
`reasoning' sub-agent that used o3-mini. The o3-mini model alone scores 87\% and
79\% on \acr{AIME} and \acr{GPQA} Diamond with a `high' reasoning effort, while
the agent system as a whole averaged 76\% across the two benchmarks.

Inspecting the traces, we observe that for many of the runs, the main agent
merely delegated the problem to the o3-mini based reasoning agent, and did not
leverage the full capabilites of the agent system. For `reasoning models' such
as o1 \citep{openaiOpenAIO1System}, o3-mini or DeepSeek's R1
\citep{deepseek-aiDeepSeekR1IncentivizingReasoning2025}), we suspect that the
inclusion of crude components aiming to induce reasoning behaviour (such as
those included in iterations 4 and 6 in Figure~\ref{fig:aime_gpqa}) may in fact
interrupt the reasoning chain of thought of a reasoning model trained outside
the agent system, resulting in a drop in performance. We look forward to future
work training `agent' \acr{LLM}s jointly with the design of the agent system
which will hopefully make better use of both components' strengths.

